% !TeX spellcheck = it_IT

\documentclass[handout]{beamer}

\usepackage[utf8]{inputenc}
\usepackage[italian]{babel}
\usepackage{multicol}
\usetheme{metropolis}

\usepackage{listings}

\usepackage{tikz}
\usetikzlibrary{shapes,arrows,positioning}

\tikzstyle{startstop} = [ellipse, draw, fill=bg, text=structure, text width=3em, text badly centered, inner sep=0pt]
\tikzstyle{block} = [rectangle, draw, fill=bg, text=structure, text width=5em, text centered]
\tikzstyle{decision} = [diamond, aspect=2, draw, fill=bg, text=structure, text width=3em, text badly centered, inner sep=0pt]
\tikzstyle{input} = [rectangle, draw, fill=bg, text=structure, text width=4em, text centered, rounded rectangle, rounded rectangle east arc=none]
\tikzstyle{output} = [rectangle, draw, fill=bg, text=structure, text width=4em, text centered, rounded rectangle, rounded rectangle west arc=none]

\tikzstyle{line} = [draw, -latex']


\title{Informatica B - Esercitazione 7}
\subtitle{Matlab}

\author{Stefano Cereda\\
		stefano1.cereda@mail.polimi.it
	}
\date{21/11/2017}
\institute[PoliMi]{\vspace{0.5cm}\centering Politecnico di Milano \\ \vspace{0.2cm}
	\includegraphics[width=0.2\linewidth]{./logopolimi}}
%\logo{\includegraphics[width=15mm]{./logopolimi}}

\setbeamercovered{invisible}

\makeindex

\begin{document}
\begin{frame}
	\maketitle
\end{frame}

\begin{frame}{Esercizio 1 - Ricerca per bisezione}
Implementare, in C ed in Matlab, l'algoritmo di ricerca dicotomica.
\end{frame}

\begin{frame}{Esercizio 2 - Aggiunta colonna}
Scrivere uno script Matlab che, data una matrice $Q$ ed un valore $v$, stampi una
matrice corrispondente a $Q$, alla quale è stata aggiunta una nuova prima
colonna contenente $v$.
\end{frame}

\begin{frame}{Esercizio 3 - Ordinamento matrice}
Porre in ordine crescente gli elementi di una matrice qualsiasi
\end{frame}

\begin{frame}{Esercizio 4  - Coppie per somme}
Scrivere uno script MATLAB che riceve in ingresso un intero positivo n e stampa due vettori x e y, contenenti tutte le coppie di interi (x,y) la cui somma è pari ad n.
\end{frame}

\begin{frame}{Esercizio 5 - Soglia intorno}
Si scriva in linguaggio MATLAB uno script che riceva in ingresso una matrice, un numero di riga, un numero di colonna e un valore di soglia. Lo script stampa true se tutti i valori delle celle della matrice nell’intorno della cella della riga e della colonna date sono minori del valore di soglia. Le celle nell’intorno della cella identificata dalla riga e dalla colonna date sono quelle immediatamente adiacenti a questa cella.
\end{frame}


\end{document}